problem:  
Let \((M^n,g)\) be a complete, noncompact Riemannian manifold with **nonnegative Ricci curvature** and **submaximal polynomial volume growth**, i.e. there exist constants \(C>0\) and \(0<\alpha<n\) such that  
\[
\operatorname{Vol}(B(p,r)) \le C\,r^{\alpha}
\quad\text{for all sufficiently large}\ r.
\]
By Gromov’s compactness theorem, for any sequence \(r_i \to \infty\), the rescaled pointed manifolds \((M, r_i^{-2}g, p)\) admit a subsequence converging in the pointed measured Gromov–Hausdorff sense to a Ricci limit space \((Y_\infty, y_\infty, \nu_\infty)\).  
Cheeger–Colding’s structure theory provides a decomposition
\[
Y_\infty = \mathcal{R} \cup \mathcal{S},
\]
where \(\mathcal{R}\) is the regular set of points with Euclidean tangent cones and \(\mathcal{S}\) is the singular set.

**Problem:**  
> Is there a universal function \(f:(0,n]\to [0,n]\) such that for every Ricci ≥ 0 manifold \((M^n,g)\) with  
> \[
> \limsup_{r\to\infty}\frac{\log \operatorname{Vol}(B(p,r))}{\log r} \le \alpha,
> \]
> **any** blow‑down Ricci limit space \(Y_\infty\) arising from \(M\) has singular set of Hausdorff dimension at most \(f(\alpha)\)?  
> 
> In particular, must one always have \(f(\alpha) < n-\alpha\)?  
>  
> Equivalently, does a quantitative upper bound on asymptotic volume growth in Ricci ≥ 0 geometry force a corresponding **quantitative upper bound on the dimension of singularities** in all asymptotic limit spaces?

---

Why is it a "good" problem:  
This problem lies squarely at the interface of **comparison geometry**, **geometric measure theory**, and **RCD limit theory**, asking whether control of global volume growth dictates the fine‑scale analytic and geometric complexity of asymptotic limit spaces.  

It corrects the logical issues of earlier formulations by:
- using \(\limsup\) instead of assuming existence of an asymptotic exponent,  
- referring to general Ricci limit spaces rather than presuming conic structure, and  
- employing the rigorous Cheeger–Colding definition of the singular set.

The question is simple yet profound: it seeks a universal law connecting **how fast volume grows** with **how badly geometry can degenerate at infinity**.  
A positive resolution would yield a new quantitative form of **asymptotic regularity**, extending classical codimension‑2 results into a dimensional family parameterized by growth rate; a counterexample would uncover unexpected flexibility in Ricci ≥ 0 limits.  

Either outcome would reshape our understanding of the interplay between global volume asymptotics and microscopic singular structures, offering an elegant “narrow entrance, profound depth” problem that could open a fresh direction in the analysis of Ricci limit spaces.